% =============================================================================
% 讲义项目级导言区(数学分析培优 / 高等代数培优 共用)
% =============================================================================
% elegantbook.cls 负责通用能力;本文件只保存讲义自己的可调设置.
% 两份讲义使用完全相同的 elegantbook.cls 与 preamble.tex,
% 页眉、水印、标题页、例题编号与数学算子因此保持一致.
% =============================================================================

% ================= 子文件支持 =================
\usepackage{subfiles}

% 单章编译可能读取 TeX Live 自带的 elegantbook.cls，故不使用文档类选项
% 控制水印，而统一读取主文件中定义的项目级开关。
\usepackage{draftwatermark}

% ================= 中文字体:楷体粗体映射 =================
% ctex 的 Windows 字体集没有给楷体声明 BoldFont,而 definition,theorem,
% remark 等环境的正文用 \cnormal (= \kaishu) 排成楷体. 于是在这些环境里
% \textbf 会去要"粗体楷体", 触发
%   LaTeX Font Warning: Font shape `TU/KaiTi(0)/b/n' undefined
% 并退回常规字重, 只靠 AutoFakeBold 合成伪粗体.
% 这里补上真粗体: 基础楷体仍是 Windows 楷体, 粗体用 SimHei, 与本书
% \heiti 已经在用的黑体是同一款, 加粗与黑体风格一致.
% 不要改写成 [BoldFont={FZHei-B01}]{FZKai-Z03}: 方括号外的 FZKai-Z03 是
% 基础字体, 那样写会连楷体字形一起换成方正楷体_GBK, 与其余仍是 Windows
% 字体的宋体,黑体,仿宋风格不统一.
\setCJKfamilyfont{zhkai}[BoldFont={SimHei}]{KaiTi}

% ================= 水印内容 =================
\SetWatermarkScale{0.4}
\SetWatermarkLightness{0.9}
\AtBeginDocument{%
    \ifmathanalysiswatermark
        \SetWatermarkText{欧拉的小迷弟}%
    \else
        \SetWatermarkText{}%
    \fi
}

% ================= 页眉页脚 =================
% ElegantBook 已加载 fancyhdr;这里只覆盖本讲义的个人化样式.
\fancyhf{}
\lhead{\rightmark}
\rhead{欧拉的小迷弟}
\cfoot{\thepage}
\renewcommand{\headrulewidth}{0.4pt}

\fancypagestyle{plain}{%
    \fancyhf{}
    \lhead{\rightmark}
    \rhead{欧拉的小迷弟}
    \cfoot{\thepage}
    \renewcommand{\headrulewidth}{0.4pt}%
}

% ================= 简洁标题页 =================
% 不使用 ElegantBook 示例封面图片,保留白底数学讲义风格.
\makeatletter
\renewcommand{\maketitle}{%
    \begin{titlepage}
        \thispagestyle{empty}
        \centering
        \vspace*{0.22\textheight}
        {\color{structurecolor}\Huge\bfseries \@title\par}
        \vspace{2.5cm}
        {\Large \@author\par}
        \vspace{1.2cm}
        {\large \@date\par}
        \vfill
    \end{titlepage}
}
\makeatother

% ================= 例题编号 =================
% 保留原讲义习惯:例题跟随 section 编号,例如“例题 1.1.1”.
\counterwithin{exam}{section}
\renewcommand{\theexam}{\thesection.\arabic{exam}}

% ================= 普通版 / 答案版留白 =================
% result=noanswer:隐藏 solution,并在例题后留 10 行;
% result=answer:显示 solution,不额外留白.
\makeatletter
\ifdefstring{\ELEGANT@result}{noanswer}{%
    \newcommand{\exampleblank}{\par\vspace{10\baselineskip}}
}{%
    \newcommand{\exampleblank}{}
}
\makeatother

% ================= 重要性标记 =================
\newcommand{\marktriangle}{\ensuremath{\blacktriangle}}
\newcommand{\markstar}{\ensuremath{\bigstar}}
\newcommand{\markstartriangle}{\ensuremath{\bigstar\blacktriangle}}
\newcommand{\marktwostar}{\ensuremath{\bigstar\bigstar}}

% ================= 章节内部的方法分类 =================
% (A)(B)(C)... 为无编号 subsubsection,同时进入目录.
% \phantomsection 用于建立正确的 PDF 超链接锚点.
\newcommand{\methodsection}[1]{%
    \subsubsection*{#1}%
    \phantomsection
    \addcontentsline{toc}{subsubsection}{#1}%
}

% ================= 通用手动留白 =================
\newcommand{\myspace}[1]{\par\vspace{#1\baselineskip}}

% ================= 自定义数学算子 =================
\DeclareMathOperator{\Tr}{Tr}
\DeclareMathOperator{\tr}{tr}
\DeclareMathOperator{\rank}{rank}
\DeclareMathOperator{\diag}{diag}
\DeclareMathOperator{\Ree}{Re}
\DeclareMathOperator{\Imm}{Im}
\DeclareMathOperator{\Ker}{Ker}
\DeclareMathOperator{\Span}{span}
% ================= 行宽溢出容错 =================
\setlength{\emergencystretch}{2em}
% ================= enumerate 列表间距 =================
\setlist[enumerate]{
    itemsep=0.4em,
    topsep=0.5em,
    parsep=0pt,
    partopsep=0pt
}
